\subsection*{Textbook Formal Inventory}
This subsection records the exact formal material in \file{Sig_ROM.ec}.  It is generated from the proof-surface EasyCrypt file, so the line numbers and statements are meant to be read next to the checked source.

\noindent\textbf{Chapter role.} generic signature, random-oracle, EUF-CMA, SUF-CMA, and UF-NMA games.

\noindent\textbf{Inventory size.} 6 context declarations, 38 definitions/game objects, 3 lemmas or relational theorems, and 0 explicit Hoare/Phoare judgments were found.

\subsubsection*{Theory Context, Imports, and Quantified Modules}
\paragraph{Context 1 (line 1)}
\noindent\textbf{EasyCrypt name.}
\begin{lstlisting}[language=EasyCrypt,basicstyle=\ttfamily\scriptsize]
imported theories
\end{lstlisting}
\noindent\textbf{Checked EasyCrypt statement.}
\begin{lstlisting}[language=EasyCrypt,basicstyle=\ttfamily\scriptsize]
require import AllCore List PROM.
\end{lstlisting}
\noindent\textbf{Formal mathematical reading.}
Mathematical context. This line imports or specializes earlier theories, making their carrier sets, functions, modules, and theorems available to the current file.

\noindent\textbf{Role in the proof.}
Use in this chapter. It fixes the proof context, imported mathematical language, or quantified module parameters for the following statements.

\paragraph{Context 2 (line 31)}
\noindent\textbf{EasyCrypt name.}
\begin{lstlisting}[language=EasyCrypt,basicstyle=\ttfamily\scriptsize]
cloned theory
\end{lstlisting}
\noindent\textbf{Checked EasyCrypt statement.}
\begin{lstlisting}[language=EasyCrypt,basicstyle=\ttfamily\scriptsize]
clone import FullRO as RO with
  type in_t <- ro_query,
  type out_t <- ro_output,
  op dout <- dro_output.
\end{lstlisting}
\noindent\textbf{Formal mathematical reading.}
Mathematical context. This line imports or specializes earlier theories, making their carrier sets, functions, modules, and theorems available to the current file.

\noindent\textbf{Role in the proof.}
Use in this chapter. It fixes the proof context, imported mathematical language, or quantified module parameters for the following statements.

\paragraph{Context 3 (line 143)}
\noindent\textbf{EasyCrypt name.}
\begin{lstlisting}[language=EasyCrypt,basicstyle=\ttfamily\scriptsize]
NMA_Embedding
\end{lstlisting}
\noindent\textbf{Checked EasyCrypt statement.}
\begin{lstlisting}[language=EasyCrypt,basicstyle=\ttfamily\scriptsize]
section NMA_Embedding.
\end{lstlisting}
\noindent\textbf{Formal mathematical reading.}
Mathematical context. This opens a scoped block of hypotheses, module parameters, and lemmas.

\noindent\textbf{Role in the proof.}
Use in this chapter. It fixes the proof context, imported mathematical language, or quantified module parameters for the following statements.

\paragraph{Context 4 (line 145)}
\noindent\textbf{EasyCrypt name.}
\begin{lstlisting}[language=EasyCrypt,basicstyle=\ttfamily\scriptsize]
H
\end{lstlisting}
\noindent\textbf{Checked EasyCrypt statement.}
\begin{lstlisting}[language=EasyCrypt,basicstyle=\ttfamily\scriptsize]
declare module H <: Oracle {-EUF_CMA}.
\end{lstlisting}
\noindent\textbf{Formal mathematical reading.}
Mathematical context. This is universal quantification over an adversary, oracle, scheme, sampler, or reduction module satisfying the stated interface and memory restrictions.

\noindent\textbf{Role in the proof.}
Use in this chapter. It fixes the proof context, imported mathematical language, or quantified module parameters for the following statements.

\paragraph{Context 5 (line 146)}
\noindent\textbf{EasyCrypt name.}
\begin{lstlisting}[language=EasyCrypt,basicstyle=\ttfamily\scriptsize]
S
\end{lstlisting}
\noindent\textbf{Checked EasyCrypt statement.}
\begin{lstlisting}[language=EasyCrypt,basicstyle=\ttfamily\scriptsize]
declare module S <: Scheme {-H, -EUF_CMA}.
\end{lstlisting}
\noindent\textbf{Formal mathematical reading.}
Mathematical context. This is universal quantification over an adversary, oracle, scheme, sampler, or reduction module satisfying the stated interface and memory restrictions.

\noindent\textbf{Role in the proof.}
Use in this chapter. It fixes the proof context, imported mathematical language, or quantified module parameters for the following statements.

\paragraph{Context 6 (line 147)}
\noindent\textbf{EasyCrypt name.}
\begin{lstlisting}[language=EasyCrypt,basicstyle=\ttfamily\scriptsize]
A
\end{lstlisting}
\noindent\textbf{Checked EasyCrypt statement.}
\begin{lstlisting}[language=EasyCrypt,basicstyle=\ttfamily\scriptsize]
declare module A <: NMA_Adversary {-H, -S, -EUF_CMA}.
\end{lstlisting}
\noindent\textbf{Formal mathematical reading.}
Mathematical context. This is universal quantification over an adversary, oracle, scheme, sampler, or reduction module satisfying the stated interface and memory restrictions.

\noindent\textbf{Role in the proof.}
Use in this chapter. It fixes the proof context, imported mathematical language, or quantified module parameters for the following statements.

\subsubsection*{Definitions, Carrier Sets, Operations, Games, and Procedures}
\begin{formaldefinition}[EasyCrypt line 5]
\noindent\textbf{EasyCrypt name.}
\begin{lstlisting}[language=EasyCrypt,basicstyle=\ttfamily\scriptsize]
pkey
\end{lstlisting}
\noindent\textbf{Checked EasyCrypt statement.}
\begin{lstlisting}[language=EasyCrypt,basicstyle=\ttfamily\scriptsize]
type pkey.
\end{lstlisting}
\noindent\textbf{Formal mathematical reading.}
Mathematical definition. \(\mathsf{pkey}\) is an abstract carrier set.  The proof may quantify over this domain, but it does not expose a representation.

\noindent\textbf{Role in the proof.}
Use in this chapter. It is part of the exact formal vocabulary used by subsequent lemmas and games.

\end{formaldefinition}

\begin{formaldefinition}[EasyCrypt line 6]
\noindent\textbf{EasyCrypt name.}
\begin{lstlisting}[language=EasyCrypt,basicstyle=\ttfamily\scriptsize]
skey
\end{lstlisting}
\noindent\textbf{Checked EasyCrypt statement.}
\begin{lstlisting}[language=EasyCrypt,basicstyle=\ttfamily\scriptsize]
type skey.
\end{lstlisting}
\noindent\textbf{Formal mathematical reading.}
Mathematical definition. \(\mathsf{skey}\) is an abstract carrier set.  The proof may quantify over this domain, but it does not expose a representation.

\noindent\textbf{Role in the proof.}
Use in this chapter. It is part of the exact formal vocabulary used by subsequent lemmas and games.

\end{formaldefinition}

\begin{formaldefinition}[EasyCrypt line 7]
\noindent\textbf{EasyCrypt name.}
\begin{lstlisting}[language=EasyCrypt,basicstyle=\ttfamily\scriptsize]
message
\end{lstlisting}
\noindent\textbf{Checked EasyCrypt statement.}
\begin{lstlisting}[language=EasyCrypt,basicstyle=\ttfamily\scriptsize]
type message.
\end{lstlisting}
\noindent\textbf{Formal mathematical reading.}
Mathematical definition. \(\mathsf{message}\) is an abstract carrier set.  The proof may quantify over this domain, but it does not expose a representation.

\noindent\textbf{Role in the proof.}
Use in this chapter. It is part of the exact formal vocabulary used by subsequent lemmas and games.

\end{formaldefinition}

\begin{formaldefinition}[EasyCrypt line 8]
\noindent\textbf{EasyCrypt name.}
\begin{lstlisting}[language=EasyCrypt,basicstyle=\ttfamily\scriptsize]
context
\end{lstlisting}
\noindent\textbf{Checked EasyCrypt statement.}
\begin{lstlisting}[language=EasyCrypt,basicstyle=\ttfamily\scriptsize]
type context.
\end{lstlisting}
\noindent\textbf{Formal mathematical reading.}
Mathematical definition. \(\mathsf{context}\) is an abstract carrier set.  The proof may quantify over this domain, but it does not expose a representation.

\noindent\textbf{Role in the proof.}
Use in this chapter. It is part of the exact formal vocabulary used by subsequent lemmas and games.

\end{formaldefinition}

\begin{formaldefinition}[EasyCrypt line 9]
\noindent\textbf{EasyCrypt name.}
\begin{lstlisting}[language=EasyCrypt,basicstyle=\ttfamily\scriptsize]
signature
\end{lstlisting}
\noindent\textbf{Checked EasyCrypt statement.}
\begin{lstlisting}[language=EasyCrypt,basicstyle=\ttfamily\scriptsize]
type signature.
\end{lstlisting}
\noindent\textbf{Formal mathematical reading.}
Mathematical definition. \(\mathsf{signature}\) is an abstract carrier set.  The proof may quantify over this domain, but it does not expose a representation.

\noindent\textbf{Role in the proof.}
Use in this chapter. It is part of the exact formal vocabulary used by subsequent lemmas and games.

\end{formaldefinition}

\begin{formaldefinition}[EasyCrypt line 10]
\noindent\textbf{EasyCrypt name.}
\begin{lstlisting}[language=EasyCrypt,basicstyle=\ttfamily\scriptsize]
ro_query
\end{lstlisting}
\noindent\textbf{Checked EasyCrypt statement.}
\begin{lstlisting}[language=EasyCrypt,basicstyle=\ttfamily\scriptsize]
type ro_query.
\end{lstlisting}
\noindent\textbf{Formal mathematical reading.}
Mathematical definition. \(\mathsf{ro\_query}\) is an abstract carrier set.  The proof may quantify over this domain, but it does not expose a representation.

\noindent\textbf{Role in the proof.}
Use in this chapter. It is part of the exact formal vocabulary used by subsequent lemmas and games.

\end{formaldefinition}

\begin{formaldefinition}[EasyCrypt line 11]
\noindent\textbf{EasyCrypt name.}
\begin{lstlisting}[language=EasyCrypt,basicstyle=\ttfamily\scriptsize]
ro_output
\end{lstlisting}
\noindent\textbf{Checked EasyCrypt statement.}
\begin{lstlisting}[language=EasyCrypt,basicstyle=\ttfamily\scriptsize]
type ro_output.
\end{lstlisting}
\noindent\textbf{Formal mathematical reading.}
Mathematical definition. \(\mathsf{ro\_output}\) is an abstract carrier set.  The proof may quantify over this domain, but it does not expose a representation.

\noindent\textbf{Role in the proof.}
Use in this chapter. It is part of the exact formal vocabulary used by subsequent lemmas and games.

\end{formaldefinition}

\begin{formaldefinition}[EasyCrypt line 13]
\noindent\textbf{EasyCrypt name.}
\begin{lstlisting}[language=EasyCrypt,basicstyle=\ttfamily\scriptsize]
dro_output
\end{lstlisting}
\noindent\textbf{Checked EasyCrypt statement.}
\begin{lstlisting}[language=EasyCrypt,basicstyle=\ttfamily\scriptsize]
op [lossless] dro_output : ro_query -> ro_output distr.
\end{lstlisting}
\noindent\textbf{Formal mathematical reading.}
Mathematical definition. This is a total EasyCrypt operation.  The exact displayed statement gives the domain, codomain, and defining equation when present.  The EasyCrypt [lossless] annotation states that this distribution has total mass 1.

\noindent\textbf{Role in the proof.}
Use in this chapter. It contributes a sampler or sampler-derived object to the distributional model.

\end{formaldefinition}

\begin{formaldefinition}[EasyCrypt line 15]
\noindent\textbf{EasyCrypt name.}
\begin{lstlisting}[language=EasyCrypt,basicstyle=\ttfamily\scriptsize]
query
\end{lstlisting}
\noindent\textbf{Checked EasyCrypt statement.}
\begin{lstlisting}[language=EasyCrypt,basicstyle=\ttfamily\scriptsize]
type query = message * context.
\end{lstlisting}
\noindent\textbf{Formal mathematical reading.}
Mathematical definition. This is a definitional type abbreviation.  The exact displayed EasyCrypt statement gives the right-hand side; later proof steps may unfold it without adding an assumption.

\noindent\textbf{Role in the proof.}
Use in this chapter. It is part of the exact formal vocabulary used by subsequent lemmas and games.

\end{formaldefinition}

\begin{formaldefinition}[EasyCrypt line 16]
\noindent\textbf{EasyCrypt name.}
\begin{lstlisting}[language=EasyCrypt,basicstyle=\ttfamily\scriptsize]
signed_query
\end{lstlisting}
\noindent\textbf{Checked EasyCrypt statement.}
\begin{lstlisting}[language=EasyCrypt,basicstyle=\ttfamily\scriptsize]
type signed_query = query * signature.
\end{lstlisting}
\noindent\textbf{Formal mathematical reading.}
Mathematical definition. This is a definitional type abbreviation.  The exact displayed EasyCrypt statement gives the right-hand side; later proof steps may unfold it without adding an assumption.

\noindent\textbf{Role in the proof.}
Use in this chapter. It is part of the exact formal vocabulary used by subsequent lemmas and games.

\end{formaldefinition}

\begin{formaldefinition}[EasyCrypt line 18]
\noindent\textbf{EasyCrypt name.}
\begin{lstlisting}[language=EasyCrypt,basicstyle=\ttfamily\scriptsize]
fresh_msg
\end{lstlisting}
\noindent\textbf{Checked EasyCrypt statement.}
\begin{lstlisting}[language=EasyCrypt,basicstyle=\ttfamily\scriptsize]
op fresh_msg (qs : query list) (m : message) (ctx : context) : bool =
  ! ((m, ctx) \in qs).
\end{lstlisting}
\noindent\textbf{Formal mathematical reading.}
Mathematical definition. This is a Boolean predicate.  The exact displayed EasyCrypt statement gives its arguments; mathematically it is an event, invariant, support condition, or well-formedness predicate over those arguments.

\noindent\textbf{Role in the proof.}
Use in this chapter. It is part of the exact formal vocabulary used by subsequent lemmas and games.

\end{formaldefinition}

\begin{formaldefinition}[EasyCrypt line 21]
\noindent\textbf{EasyCrypt name.}
\begin{lstlisting}[language=EasyCrypt,basicstyle=\ttfamily\scriptsize]
fresh_sig
\end{lstlisting}
\noindent\textbf{Checked EasyCrypt statement.}
\begin{lstlisting}[language=EasyCrypt,basicstyle=\ttfamily\scriptsize]
op fresh_sig (qs : signed_query list) (m : message) (ctx : context)
             (sig : signature) : bool =
  ! (((m, ctx), sig) \in qs).
\end{lstlisting}
\noindent\textbf{Formal mathematical reading.}
Mathematical definition. This is a Boolean predicate.  The exact displayed EasyCrypt statement gives its arguments; mathematically it is an event, invariant, support condition, or well-formedness predicate over those arguments.

\noindent\textbf{Role in the proof.}
Use in this chapter. It is part of the exact formal vocabulary used by subsequent lemmas and games.

\end{formaldefinition}

\begin{formaldefinition}[EasyCrypt line 36]
\noindent\textbf{EasyCrypt name.}
\begin{lstlisting}[language=EasyCrypt,basicstyle=\ttfamily\scriptsize]
Oracle
\end{lstlisting}
\noindent\textbf{Checked EasyCrypt statement.}
\begin{lstlisting}[language=EasyCrypt,basicstyle=\ttfamily\scriptsize]
module type Oracle = {
\end{lstlisting}
\noindent\textbf{Formal mathematical reading.}
Mathematical definition. This is an interface for an oracle, adversary, scheme, sampler, solver, or reduction.  A later theorem quantified over this interface is universally quantified over all modules implementing these procedures.

\noindent\textbf{Role in the proof.}
Use in this chapter. It defines the oracle boundary through which adversaries and simulations interact.

\end{formaldefinition}

\begin{formaldefinition}[EasyCrypt line 40]
\noindent\textbf{EasyCrypt name.}
\begin{lstlisting}[language=EasyCrypt,basicstyle=\ttfamily\scriptsize]
POracle
\end{lstlisting}
\noindent\textbf{Checked EasyCrypt statement.}
\begin{lstlisting}[language=EasyCrypt,basicstyle=\ttfamily\scriptsize]
module type POracle = {
\end{lstlisting}
\noindent\textbf{Formal mathematical reading.}
Mathematical definition. This is an interface for an oracle, adversary, scheme, sampler, solver, or reduction.  A later theorem quantified over this interface is universally quantified over all modules implementing these procedures.

\noindent\textbf{Role in the proof.}
Use in this chapter. It defines the oracle boundary through which adversaries and simulations interact.

\end{formaldefinition}

\begin{formaldefinition}[EasyCrypt line 44]
\noindent\textbf{EasyCrypt name.}
\begin{lstlisting}[language=EasyCrypt,basicstyle=\ttfamily\scriptsize]
Scheme
\end{lstlisting}
\noindent\textbf{Checked EasyCrypt statement.}
\begin{lstlisting}[language=EasyCrypt,basicstyle=\ttfamily\scriptsize]
module type Scheme(O : POracle) = {
\end{lstlisting}
\noindent\textbf{Formal mathematical reading.}
Mathematical definition. This is an interface for an oracle, adversary, scheme, sampler, solver, or reduction.  A later theorem quantified over this interface is universally quantified over all modules implementing these procedures.

\noindent\textbf{Role in the proof.}
Use in this chapter. It supplies an executable component of the formal probabilistic model.

\end{formaldefinition}

\begin{formaldefinition}[EasyCrypt line 45]
\noindent\textbf{EasyCrypt name.}
\begin{lstlisting}[language=EasyCrypt,basicstyle=\ttfamily\scriptsize]
kg
\end{lstlisting}
\noindent\textbf{Checked EasyCrypt statement.}
\begin{lstlisting}[language=EasyCrypt,basicstyle=\ttfamily\scriptsize]
proc kg() : pkey * skey
  proc sign(sk : skey, m : message, ctx : context) : signature
  proc verify(pk : pkey, m : message, ctx : context,
              sig : signature) : bool
}.
\end{lstlisting}
\noindent\textbf{Formal mathematical reading.}
Mathematical definition. This procedure is a command in the probabilistic program.  Its return value, state updates, and oracle calls are the operational content of the game.

\noindent\textbf{Role in the proof.}
Use in this chapter. It supplies an executable component of the formal probabilistic model.

\end{formaldefinition}

\begin{formaldefinition}[EasyCrypt line 51]
\noindent\textbf{EasyCrypt name.}
\begin{lstlisting}[language=EasyCrypt,basicstyle=\ttfamily\scriptsize]
CORR_ADV
\end{lstlisting}
\noindent\textbf{Checked EasyCrypt statement.}
\begin{lstlisting}[language=EasyCrypt,basicstyle=\ttfamily\scriptsize]
module type CORR_ADV(O : POracle) = {
\end{lstlisting}
\noindent\textbf{Formal mathematical reading.}
Mathematical definition. This is an interface for an oracle, adversary, scheme, sampler, solver, or reduction.  A later theorem quantified over this interface is universally quantified over all modules implementing these procedures.

\noindent\textbf{Role in the proof.}
Use in this chapter. It supplies an executable component of the formal probabilistic model.

\end{formaldefinition}

\begin{formaldefinition}[EasyCrypt line 52]
\noindent\textbf{EasyCrypt name.}
\begin{lstlisting}[language=EasyCrypt,basicstyle=\ttfamily\scriptsize]
find
\end{lstlisting}
\noindent\textbf{Checked EasyCrypt statement.}
\begin{lstlisting}[language=EasyCrypt,basicstyle=\ttfamily\scriptsize]
proc find(pk : pkey, sk : skey) : message * context
}.
\end{lstlisting}
\noindent\textbf{Formal mathematical reading.}
Mathematical definition. This procedure is a command in the probabilistic program.  Its return value, state updates, and oracle calls are the operational content of the game.

\noindent\textbf{Role in the proof.}
Use in this chapter. It supplies an executable component of the formal probabilistic model.

\end{formaldefinition}

\begin{formaldefinition}[EasyCrypt line 55]
\noindent\textbf{EasyCrypt name.}
\begin{lstlisting}[language=EasyCrypt,basicstyle=\ttfamily\scriptsize]
Correctness
\end{lstlisting}
\noindent\textbf{Checked EasyCrypt statement.}
\begin{lstlisting}[language=EasyCrypt,basicstyle=\ttfamily\scriptsize]
module Correctness(H : Oracle, S : Scheme, A : CORR_ADV) = {
\end{lstlisting}
\noindent\textbf{Formal mathematical reading.}
Mathematical definition. This is a probabilistic program/game or a module adapter.  Its procedures induce the probability space used by the game-based proof.

\noindent\textbf{Role in the proof.}
Use in this chapter. It supplies an executable component of the formal probabilistic model.

\end{formaldefinition}

\begin{formaldefinition}[EasyCrypt line 56]
\noindent\textbf{EasyCrypt name.}
\begin{lstlisting}[language=EasyCrypt,basicstyle=\ttfamily\scriptsize]
main
\end{lstlisting}
\noindent\textbf{Checked EasyCrypt statement.}
\begin{lstlisting}[language=EasyCrypt,basicstyle=\ttfamily\scriptsize]
proc main() : bool = {
\end{lstlisting}
\noindent\textbf{Formal mathematical reading.}
Mathematical definition. This procedure is a command in the probabilistic program.  Its return value, state updates, and oracle calls are the operational content of the game.

\noindent\textbf{Role in the proof.}
Use in this chapter. It supplies an executable component of the formal probabilistic model.

\end{formaldefinition}

\begin{formaldefinition}[EasyCrypt line 73]
\noindent\textbf{EasyCrypt name.}
\begin{lstlisting}[language=EasyCrypt,basicstyle=\ttfamily\scriptsize]
SignOracle
\end{lstlisting}
\noindent\textbf{Checked EasyCrypt statement.}
\begin{lstlisting}[language=EasyCrypt,basicstyle=\ttfamily\scriptsize]
module type SignOracle = {
\end{lstlisting}
\noindent\textbf{Formal mathematical reading.}
Mathematical definition. This is an interface for an oracle, adversary, scheme, sampler, solver, or reduction.  A later theorem quantified over this interface is universally quantified over all modules implementing these procedures.

\noindent\textbf{Role in the proof.}
Use in this chapter. It defines the oracle boundary through which adversaries and simulations interact.

\end{formaldefinition}

\begin{formaldefinition}[EasyCrypt line 74]
\noindent\textbf{EasyCrypt name.}
\begin{lstlisting}[language=EasyCrypt,basicstyle=\ttfamily\scriptsize]
sign
\end{lstlisting}
\noindent\textbf{Checked EasyCrypt statement.}
\begin{lstlisting}[language=EasyCrypt,basicstyle=\ttfamily\scriptsize]
proc sign(m : message, ctx : context) : signature
}.
\end{lstlisting}
\noindent\textbf{Formal mathematical reading.}
Mathematical definition. This procedure is a command in the probabilistic program.  Its return value, state updates, and oracle calls are the operational content of the game.

\noindent\textbf{Role in the proof.}
Use in this chapter. It supplies an executable component of the formal probabilistic model.

\end{formaldefinition}

\begin{formaldefinition}[EasyCrypt line 77]
\noindent\textbf{EasyCrypt name.}
\begin{lstlisting}[language=EasyCrypt,basicstyle=\ttfamily\scriptsize]
Adversary
\end{lstlisting}
\noindent\textbf{Checked EasyCrypt statement.}
\begin{lstlisting}[language=EasyCrypt,basicstyle=\ttfamily\scriptsize]
module type Adversary(H : POracle, O : SignOracle) = {
\end{lstlisting}
\noindent\textbf{Formal mathematical reading.}
Mathematical definition. This is an interface for an oracle, adversary, scheme, sampler, solver, or reduction.  A later theorem quantified over this interface is universally quantified over all modules implementing these procedures.

\noindent\textbf{Role in the proof.}
Use in this chapter. It supplies an executable component of the formal probabilistic model.

\end{formaldefinition}

\begin{formaldefinition}[EasyCrypt line 78]
\noindent\textbf{EasyCrypt name.}
\begin{lstlisting}[language=EasyCrypt,basicstyle=\ttfamily\scriptsize]
forge
\end{lstlisting}
\noindent\textbf{Checked EasyCrypt statement.}
\begin{lstlisting}[language=EasyCrypt,basicstyle=\ttfamily\scriptsize]
proc forge(pk : pkey) : message * context * signature {H.get, O.sign}
\end{lstlisting}
\noindent\textbf{Formal mathematical reading.}
Mathematical definition. This procedure is a command in the probabilistic program.  Its return value, state updates, and oracle calls are the operational content of the game.

\noindent\textbf{Role in the proof.}
Use in this chapter. It supplies an executable component of the formal probabilistic model.

\end{formaldefinition}

\begin{formaldefinition}[EasyCrypt line 81]
\noindent\textbf{EasyCrypt name.}
\begin{lstlisting}[language=EasyCrypt,basicstyle=\ttfamily\scriptsize]
NMA_Adversary
\end{lstlisting}
\noindent\textbf{Checked EasyCrypt statement.}
\begin{lstlisting}[language=EasyCrypt,basicstyle=\ttfamily\scriptsize]
module type NMA_Adversary(H : POracle) = {
\end{lstlisting}
\noindent\textbf{Formal mathematical reading.}
Mathematical definition. This is an interface for an oracle, adversary, scheme, sampler, solver, or reduction.  A later theorem quantified over this interface is universally quantified over all modules implementing these procedures.

\noindent\textbf{Role in the proof.}
Use in this chapter. It defines the game or game interface whose success event is bounded by later theorems.

\end{formaldefinition}

\begin{formaldefinition}[EasyCrypt line 82]
\noindent\textbf{EasyCrypt name.}
\begin{lstlisting}[language=EasyCrypt,basicstyle=\ttfamily\scriptsize]
forge
\end{lstlisting}
\noindent\textbf{Checked EasyCrypt statement.}
\begin{lstlisting}[language=EasyCrypt,basicstyle=\ttfamily\scriptsize]
proc forge(pk : pkey) : message * context * signature {H.get}
\end{lstlisting}
\noindent\textbf{Formal mathematical reading.}
Mathematical definition. This procedure is a command in the probabilistic program.  Its return value, state updates, and oracle calls are the operational content of the game.

\noindent\textbf{Role in the proof.}
Use in this chapter. It supplies an executable component of the formal probabilistic model.

\end{formaldefinition}

\begin{formaldefinition}[EasyCrypt line 85]
\noindent\textbf{EasyCrypt name.}
\begin{lstlisting}[language=EasyCrypt,basicstyle=\ttfamily\scriptsize]
NMA_As_CMA
\end{lstlisting}
\noindent\textbf{Checked EasyCrypt statement.}
\begin{lstlisting}[language=EasyCrypt,basicstyle=\ttfamily\scriptsize]
module NMA_As_CMA(A : NMA_Adversary) (H : POracle, O : SignOracle) = {
\end{lstlisting}
\noindent\textbf{Formal mathematical reading.}
Mathematical definition. This is a probabilistic program/game or a module adapter.  Its procedures induce the probability space used by the game-based proof.

\noindent\textbf{Role in the proof.}
Use in this chapter. It defines the game or game interface whose success event is bounded by later theorems.

\end{formaldefinition}

\begin{formaldefinition}[EasyCrypt line 86]
\noindent\textbf{EasyCrypt name.}
\begin{lstlisting}[language=EasyCrypt,basicstyle=\ttfamily\scriptsize]
forge
\end{lstlisting}
\noindent\textbf{Checked EasyCrypt statement.}
\begin{lstlisting}[language=EasyCrypt,basicstyle=\ttfamily\scriptsize]
proc forge(pk : pkey) : message * context * signature = {
\end{lstlisting}
\noindent\textbf{Formal mathematical reading.}
Mathematical definition. This procedure is a command in the probabilistic program.  Its return value, state updates, and oracle calls are the operational content of the game.

\noindent\textbf{Role in the proof.}
Use in this chapter. It supplies an executable component of the formal probabilistic model.

\end{formaldefinition}

\begin{formaldefinition}[EasyCrypt line 94]
\noindent\textbf{EasyCrypt name.}
\begin{lstlisting}[language=EasyCrypt,basicstyle=\ttfamily\scriptsize]
EUF_CMA
\end{lstlisting}
\noindent\textbf{Checked EasyCrypt statement.}
\begin{lstlisting}[language=EasyCrypt,basicstyle=\ttfamily\scriptsize]
module EUF_CMA(H : Oracle, S : Scheme, A : Adversary) = {
\end{lstlisting}
\noindent\textbf{Formal mathematical reading.}
Mathematical definition. This is a probabilistic program/game or a module adapter.  Its procedures induce the probability space used by the game-based proof.

\noindent\textbf{Role in the proof.}
Use in this chapter. It defines the game or game interface whose success event is bounded by later theorems.

\end{formaldefinition}

\begin{formaldefinition}[EasyCrypt line 98]
\noindent\textbf{EasyCrypt name.}
\begin{lstlisting}[language=EasyCrypt,basicstyle=\ttfamily\scriptsize]
O
\end{lstlisting}
\noindent\textbf{Checked EasyCrypt statement.}
\begin{lstlisting}[language=EasyCrypt,basicstyle=\ttfamily\scriptsize]
module O = {
\end{lstlisting}
\noindent\textbf{Formal mathematical reading.}
Mathematical definition. This is a probabilistic program/game or a module adapter.  Its procedures induce the probability space used by the game-based proof.

\noindent\textbf{Role in the proof.}
Use in this chapter. It supplies an executable component of the formal probabilistic model.

\end{formaldefinition}

\begin{formaldefinition}[EasyCrypt line 99]
\noindent\textbf{EasyCrypt name.}
\begin{lstlisting}[language=EasyCrypt,basicstyle=\ttfamily\scriptsize]
sign
\end{lstlisting}
\noindent\textbf{Checked EasyCrypt statement.}
\begin{lstlisting}[language=EasyCrypt,basicstyle=\ttfamily\scriptsize]
proc sign(m : message, ctx : context) : signature = {
\end{lstlisting}
\noindent\textbf{Formal mathematical reading.}
Mathematical definition. This procedure is a command in the probabilistic program.  Its return value, state updates, and oracle calls are the operational content of the game.

\noindent\textbf{Role in the proof.}
Use in this chapter. It supplies an executable component of the formal probabilistic model.

\end{formaldefinition}

\begin{formaldefinition}[EasyCrypt line 108]
\noindent\textbf{EasyCrypt name.}
\begin{lstlisting}[language=EasyCrypt,basicstyle=\ttfamily\scriptsize]
A
\end{lstlisting}
\noindent\textbf{Checked EasyCrypt statement.}
\begin{lstlisting}[language=EasyCrypt,basicstyle=\ttfamily\scriptsize]
module A = A(H, O)

  proc main() : bool = {
\end{lstlisting}
\noindent\textbf{Formal mathematical reading.}
Mathematical definition. This is a probabilistic program/game or a module adapter.  Its procedures induce the probability space used by the game-based proof.

\noindent\textbf{Role in the proof.}
Use in this chapter. It supplies an executable component of the formal probabilistic model.

\end{formaldefinition}

\begin{formaldefinition}[EasyCrypt line 126]
\noindent\textbf{EasyCrypt name.}
\begin{lstlisting}[language=EasyCrypt,basicstyle=\ttfamily\scriptsize]
UF_NMA
\end{lstlisting}
\noindent\textbf{Checked EasyCrypt statement.}
\begin{lstlisting}[language=EasyCrypt,basicstyle=\ttfamily\scriptsize]
module UF_NMA(H : Oracle, S : Scheme, A : NMA_Adversary) = {
\end{lstlisting}
\noindent\textbf{Formal mathematical reading.}
Mathematical definition. This is a probabilistic program/game or a module adapter.  Its procedures induce the probability space used by the game-based proof.

\noindent\textbf{Role in the proof.}
Use in this chapter. It defines the game or game interface whose success event is bounded by later theorems.

\end{formaldefinition}

\begin{formaldefinition}[EasyCrypt line 127]
\noindent\textbf{EasyCrypt name.}
\begin{lstlisting}[language=EasyCrypt,basicstyle=\ttfamily\scriptsize]
main
\end{lstlisting}
\noindent\textbf{Checked EasyCrypt statement.}
\begin{lstlisting}[language=EasyCrypt,basicstyle=\ttfamily\scriptsize]
proc main() : bool = {
\end{lstlisting}
\noindent\textbf{Formal mathematical reading.}
Mathematical definition. This procedure is a command in the probabilistic program.  Its return value, state updates, and oracle calls are the operational content of the game.

\noindent\textbf{Role in the proof.}
Use in this chapter. It supplies an executable component of the formal probabilistic model.

\end{formaldefinition}

\begin{formaldefinition}[EasyCrypt line 174]
\noindent\textbf{EasyCrypt name.}
\begin{lstlisting}[language=EasyCrypt,basicstyle=\ttfamily\scriptsize]
SUF_CMA
\end{lstlisting}
\noindent\textbf{Checked EasyCrypt statement.}
\begin{lstlisting}[language=EasyCrypt,basicstyle=\ttfamily\scriptsize]
module SUF_CMA(H : Oracle, S : Scheme, A : Adversary) = {
\end{lstlisting}
\noindent\textbf{Formal mathematical reading.}
Mathematical definition. This is a probabilistic program/game or a module adapter.  Its procedures induce the probability space used by the game-based proof.

\noindent\textbf{Role in the proof.}
Use in this chapter. It defines the game or game interface whose success event is bounded by later theorems.

\end{formaldefinition}

\begin{formaldefinition}[EasyCrypt line 178]
\noindent\textbf{EasyCrypt name.}
\begin{lstlisting}[language=EasyCrypt,basicstyle=\ttfamily\scriptsize]
O
\end{lstlisting}
\noindent\textbf{Checked EasyCrypt statement.}
\begin{lstlisting}[language=EasyCrypt,basicstyle=\ttfamily\scriptsize]
module O = {
\end{lstlisting}
\noindent\textbf{Formal mathematical reading.}
Mathematical definition. This is a probabilistic program/game or a module adapter.  Its procedures induce the probability space used by the game-based proof.

\noindent\textbf{Role in the proof.}
Use in this chapter. It supplies an executable component of the formal probabilistic model.

\end{formaldefinition}

\begin{formaldefinition}[EasyCrypt line 179]
\noindent\textbf{EasyCrypt name.}
\begin{lstlisting}[language=EasyCrypt,basicstyle=\ttfamily\scriptsize]
sign
\end{lstlisting}
\noindent\textbf{Checked EasyCrypt statement.}
\begin{lstlisting}[language=EasyCrypt,basicstyle=\ttfamily\scriptsize]
proc sign(m : message, ctx : context) : signature = {
\end{lstlisting}
\noindent\textbf{Formal mathematical reading.}
Mathematical definition. This procedure is a command in the probabilistic program.  Its return value, state updates, and oracle calls are the operational content of the game.

\noindent\textbf{Role in the proof.}
Use in this chapter. It supplies an executable component of the formal probabilistic model.

\end{formaldefinition}

\begin{formaldefinition}[EasyCrypt line 188]
\noindent\textbf{EasyCrypt name.}
\begin{lstlisting}[language=EasyCrypt,basicstyle=\ttfamily\scriptsize]
A
\end{lstlisting}
\noindent\textbf{Checked EasyCrypt statement.}
\begin{lstlisting}[language=EasyCrypt,basicstyle=\ttfamily\scriptsize]
module A = A(H, O)

  proc main() : bool = {
\end{lstlisting}
\noindent\textbf{Formal mathematical reading.}
Mathematical definition. This is a probabilistic program/game or a module adapter.  Its procedures induce the probability space used by the game-based proof.

\noindent\textbf{Role in the proof.}
Use in this chapter. It supplies an executable component of the formal probabilistic model.

\end{formaldefinition}

\subsubsection*{Lemmas, Theorems, Relational Equivalences, and Their Proof Dependencies}
\begin{formaltheorem}[EasyCrypt line 25]
\noindent\textbf{EasyCrypt name.}
\begin{lstlisting}[language=EasyCrypt,basicstyle=\ttfamily\scriptsize]
fresh_msg_nil
\end{lstlisting}
\noindent\textbf{Checked EasyCrypt statement.}
\begin{lstlisting}[language=EasyCrypt,basicstyle=\ttfamily\scriptsize]
lemma fresh_msg_nil m ctx : fresh_msg [] m ctx.
\end{lstlisting}
\noindent\textbf{Formal mathematical reading.}
Mathematical theorem. This lemma is a universally quantified proposition over the variables shown in the EasyCrypt statement.

\noindent\textbf{Role in the proof.}
Use in this chapter. It is a reusable checked theorem cited by later proof scripts.

\noindent\textbf{Proof-script interpretation.}
No proof block was collected for this item.

\end{formaltheorem}

\begin{formaltheorem}[EasyCrypt line 28]
\noindent\textbf{EasyCrypt name.}
\begin{lstlisting}[language=EasyCrypt,basicstyle=\ttfamily\scriptsize]
fresh_sig_nil
\end{lstlisting}
\noindent\textbf{Checked EasyCrypt statement.}
\begin{lstlisting}[language=EasyCrypt,basicstyle=\ttfamily\scriptsize]
lemma fresh_sig_nil m ctx sig : fresh_sig [] m ctx sig.
\end{lstlisting}
\noindent\textbf{Formal mathematical reading.}
Mathematical theorem. This lemma is a universally quantified proposition over the variables shown in the EasyCrypt statement.

\noindent\textbf{Role in the proof.}
Use in this chapter. It is a reusable checked theorem cited by later proof scripts.

\noindent\textbf{Proof-script interpretation.}
No proof block was collected for this item.

\end{formaltheorem}

\begin{formaltheorem}[EasyCrypt line 149]
\noindent\textbf{EasyCrypt name.}
\begin{lstlisting}[language=EasyCrypt,basicstyle=\ttfamily\scriptsize]
nma_as_cma_exact
\end{lstlisting}
\noindent\textbf{Checked EasyCrypt statement.}
\begin{lstlisting}[language=EasyCrypt,basicstyle=\ttfamily\scriptsize]
lemma nma_as_cma_exact &m :
  Pr[EUF_CMA(H, S, NMA_As_CMA(A)).main() @ &m : res] =
  Pr[UF_NMA(H, S, A).main() @ &m : res].
\end{lstlisting}
\noindent\textbf{Formal mathematical reading.}
Mathematical theorem. This theorem has the form \(\Pr[G:E] = B\) is a game-probability statement.  It is one of the exact hops, bad-event bounds, or final advantage inequalities used in the reduction.

\noindent\textbf{Role in the proof.}
Use in this chapter. It proves an exact game replacement or simulator equivalence.

\noindent\textbf{Proof-script interpretation.}
Proof-script reading. The machine proof decomposes the theorem into rewrites, program-logic obligations, relational equivalences, and arithmetic side conditions.
\emph{Main tactics visible in the proof script:} \ec{byequiv}, \ec{proc}, \ec{inline}, \ec{wp}, \ec{call}.
\emph{Named identifiers syntactically cited in the proof block:}
\begin{lstlisting}[language=EasyCrypt,basicstyle=\ttfamily\scriptsize]
EUF_CMA
NMA_As_CMA
arg
forge
glob
\end{lstlisting}

\end{formaltheorem}

\medskip
\noindent\emph{End of generated formal inventory for \file{Sig_ROM.ec}.}
